\section{Обзор работ по прогнозированию динамики графов транзакций}

\subsection*{Резюме} 
Прогнозирование динамики графов транзакций рассматривается в различных контекстах: от общего моделирования эволюции графов до специфичных задач в криптовалютах и платежных сетях. В работах изучаются такие цели, как предсказание появления новых ребер (link prediction), прогнозирование весов/объёмов транзакций (flow prediction), прогнозирование временных моментов и интервалов транзакций, обнаружение аномалий (мошенничества) и прогнозирование свойств узлов. Графы могут моделироваться как статические или динамические (либо дискретно-временные последовательности, либо потоки событий). Методы варьируются от статистических и классических моделей временных рядов до нейросетевых подходов: графовых нейронных сетей (динамических GNN), временных моделях (RNN/Transformer), либо указанных графовых эмбеддингов. Для каждой категории рассмотрены ключевые подходы и примеры статей. В частности, работы [1]–[6], [8], [24], [26], [33], [41] демонстрируют использование гибридных подходов: например, предсказание структуры будущего графа на основе временных рядов степеней узлов и задачи оптимизации【1†L43-L51】, использование графовых нейронных сетей с учётом временных меток и весов ребер【6†L22-L30】【8†L39-L48】, а также специализированные решения для финансовых сетей (мультиграфы, весовые атрибуты)【8†L113-L121】【41†L49-L56】. Масштаб данных (миллионы узлов/ребер) и высокая активность транзакций требовали дискретизации во времени (например, дневные срезы) для повышения качества прогноза【17†L859-L863】【15†L731-L739】. Из анализа литературы видно, что link prediction является естественной постановкой для прогноза транзакций, поскольку цель (новые транзакции) доступны из данных, а метрики (AUC, MRR) хорошо подходят для оценки. По сравнению с альтернативами (регрессия весов, прогноз суммарных метрик, обнаружение аномалий) link prediction обеспечивает простоту меток и коррелирует напрямую с задачей прогнозирования сетевой активности. Ниже приводится классификация и сравнение ключевых работ по указанным критериям, заключения и аргументы в пользу link prediction для нашей задачи.

\subsection{Постановки задач прогнозирования}
Среди главных исследуемых задач выделяются:
\begin{itemize}
\item \textbf{Link prediction (прогноз ребер)}: предсказание появления новых транзакционных связей между узлами. Примеры: Wei et al. (2022) строят прогноз транзакций между Bitcoin-аккаунтами с помощью глубокой сети, обученной на временных графах【24†L9-L17】; Tang et al. (2023) исследуют временное link prediction в NFT-сети Ethereum, используя GNN на ежедневных срезах с метриками AUC/MRR【17†L859-L863】【15†L731-L739】.
\item \textbf{Flow prediction (прогноз весов)}: прогноз объёмов/весов ребер (существующих или новых). Напр., Takahashi et al. (2024) одновременно прогнозируют наличие ребра и объём платежа, разделяя задачу на прогноз соотношения переводов и общего объёма, используя механизм само-внимания【6†L22-L30】. Liang et al. (2024) предлагают LineGNN, преобразующий локальный субграф в ребро-граф и обучающий GCN для предсказания веса заданного ребра【22†L9-L17】. Эти задачи обычно формулируются как регрессия на известных связях【6†L22-L30】【22†L9-L17】.
\item \textbf{Temporal prediction (прогноз времен)}: предсказание моментов/интервалов транзакций либо поведения активности узлов. Это направление тесно связано с подходами временных точечных процессов (TPP) или моделями очереди, но в арXив-литературе специфические работы по прогнозу времён транзакций малочисленны. Скорее, прогноз временных показателей (например, средняя задержка, частота транзакций) рассматривается через временные ряды [26].
\item \textbf{Anomaly detection (обнаружение аномалий/мошенничество)}: поиск мошеннических узлов или транзакций. Здесь задача обычно формулируется как классификация узлов/ребер. Например, финансовые графы транзакций анализируются GNN для AML-задач【8†L39-L48】【41†L49-L56】. Недавние обзоры подчёркивают эффективность GNN в захвате сложных связей и динамики сети, существенно превосходя традиционные методы для обнаружения мошенничества【19†L60-L63】.
\item \textbf{Node property prediction (свойства узлов)}: прогнозирования атрибутов узлов, например категорий участников, склонности к активности. Ghimire et al. (2026) сравнили гиперболические и евклидовы GNN для классификации адресов Bitcoin и обнаружили, что гиперболические эмбеддинги дают лучшие результаты при больших глубинах сети【33†L166-L174】. Это важно для задач сущностной классификации (классификация типов аккаунтов, рисковых узлов и т.п.).
\end{itemize}

\subsection{Временная модель графа}
Работы по динамическим графам транзакций используют разные модели времени. В дискретно-временной модели граф аппроксимируется последовательностью графов-срезов (ежедневных, еженедельных, ежемесячных и т.д.)【15†L731-L739】【17†L859-L863】. Например, Tang et al. (2023) формируют серию ежедневных графов из данных Ethereum-NFT и учат GNN на предсказание ребер в следующем срезе【17†L859-L863】. Другой подход — непрерывно-временные динамические графы, где события транзакций рассматриваются как поток событий (continuous-time dynamic graphs). Для их моделирования применяют temporal GNN (TGAT, TGN) или рекуррентные сети над событиями. Некоторые методы (например, Temporal Graph Attention) учитывают таймстемпы непосредственно в агрегации. Kandanaarachchi (2024) рассматривает дискретные динамические графы и даже допускает появление новых узлов при прогнозировании будущего графа【1†L43-L51】. В целом, выбор модели (статический vs динамический, дискретный vs непрерывный) зависит от плотности данных и задач: при высокой активности транзакций дискретизация по малым интервалам (сутки) улучшает качество прогноза, как показано для NFT-графа【17†L859-L863】【15†L731-L739】. 

\subsection{Методы и подходы}
Подходы к прогнозированию в графах транзакций включают:
\begin{itemize}
\item \textbf{Статистические и классические модели}. Например, в задаче прогнозирования сборов транзакций SARIMAX и Prophet превосходят сложные нейросетевые модели при ограниченных данных【26†L61-L69】. Time-series методы (SARIMA, Prophet, Hybrid) хорошо подходят для предсказания агрегированных метрик (fees, объёмы), но плохо масштабируются на детали структуры графа.
\item \textbf{Графовые нейронные сети (GNN)}. Для link prediction в динамических графах применяют Temporal GNN, комбинирующие топологию и время. Takahashi et al. используют GNN с временным само-вниманием (TGAT-подобный слой) для прогноза напралений потоков【6†L22-L30】. TeMP-TraG (2025) вводит временной компонент в message passing, взвешивая недавние транзакции выше【8†L113-L121】. GNN-архитектуры включают EvolveGCN, DynGEM, GRU/GAT-комбинации и др. Tang et al. сравнивают несколько таких моделей (GCRN, Roland, DynGEM) на временных графах Ethereum и показывают, что специально обученный метод Roland даёт лучшие результаты【17†L859-L863】. 
\item \textbf{Последовательные модели и представления}. Некоторые работы используют рекуррентные или трансформер-модели поверх графовых эмбеддингов. Wei et al. (2022) строят две сети на основе эмбеддингов узлов: одна для доступности (reachability), другая — для динамики сумм, и ансамблируют их【24†L15-L24】. Это позволяет захватить нелинейные пространственно-временные паттерны в Bitcoin-графе.
\item \textbf{Временные точечные процессы (TPP)}. Подобные методы (Hawkes-процессы, neural TPP) рассматривают транзакции как события. Хотя прямых арxiv-работ по криптографическим данным немного, в целом TPP могут моделировать зависимость интервала событий от истории, что полезно для прогнозирования момента следующей транзакции (future transaction time). ТPP часто применяют в сетевых потоках информации, и их можно адаптировать к финансовым потокам.
\item \textbf{Методы представления (Representation learning)}. Включают статические эмбеддинги (node2vec, DeepWalk) и глубокие представления ребер (LineGNN)【22†L9-L17】. Лианг и др. применяют преобразование в line-граф, чтобы обучить CNN для задачи регрессии веса ребра, показывая улучшение над классическими методами【22†L9-L17】. 
\end{itemize}

\subsection{Метрики оценки}
Для link prediction обычно используют \emph{AUC}, \emph{ROC-AUC}, \emph{MRR} и точностные метрики (Precision@k)【15†L731-L739】【17†L859-L863】. Для предсказания весов – среднюю абсолютную ошибку (MAE), MAPE, RMSE. Для классификации узлов/ребер (аномалий) – точность, полноту, F-мера【19†L60-L63】【33†L166-L174】. В обзорах отмечается, что в задачах динамики важно корректное разбиение на обучение/тест (fixed split vs rolling) для адекватной оценки【15†L731-L739】. 

\subsection{Ключевые примеры и группировка работ}
Ниже мы группируем работы по типам задач и приводим ключевые примеры с кратким описанием:

\paragraph{Link prediction в транзакционных графах.}  Wei et al. (2022) решают задачу прогнозирования транзакций Bitcoin-сети, строя \emph{time-decaying reachability} и \emph{amount graphs}, обучая эмбеддинги узлов и используя их для вероятностной оценки будущих ребер【24†L15-L24】. Они достигают ~60\% точности прогнозирования и выигрыша в 1.5 раза перед статической моделью. Tang et al. (2023) анализируют NFT-граф Ethereum: они используют GNN (Dyngraph2Vec, GCRN, а также мета-обученную \emph{Roland}) для прогнозирования будущих взаимодействий, и находят, что более мелкая дискретизация (день vs месяц) даёт значительно лучшие AUC【17†L859-L863】【15†L731-L739】, а подходы типа Roland превосходят базовые (особенно версия с moving average)【17†L859-L863】. Kandanaarachchi (2024) предлагает комбинировать прогноз степеней узлов (TS-прогноз) и оптимизационный подход Flux Balance Analysis для предсказания целой структуры будущего графа【1†L43-L51】. Эта работа демонстрирует новую постановку – предсказывать все узлы и ребра в будущем шаге, а не только скрытые ребра в текущем графе. EX-Graph (Wang et al., 2023) предоставляет большой датасет Ethereum+социальные связи с задачей link prediction (и детектирования wash-трейдинга)【38†L1-L4】, подчёркивая важность объединения разных источников для предсказания связей. 

\paragraph{Flow/weight prediction.} Takahashi et al. (2024) одновременно прогнозируют наличие ребра и его вес: они разделяют задачу на прогноз отношения потоков (remittance ratios) и общего объёма по узлам【6†L22-L30】. Модель использует временное само-внимание, объединяя топологию и время. В экспериментах на криптовалютном и банковском данных такой подход показывает высокую точность. Liang et al. (2024) декларируют, что для весов ребер важно обучать специфические представления звеньев: их LGLWP-технология сначала извлекает окрестность ребра, строит line-граф и применяет GNN, после чего регрессия даёт вес. Это прямолинейно учит реберные признаки и превосходит предыдущие методы по MAE【22†L9-L17】. Отметим, что в flow prediction обычно фиксируются существующие связи (динамика весов на них), тогда как link prediction нацелен именно на появление новых связей【6†L22-L30】【22†L9-L17】. 

\paragraph{Anomaly detection (AML, fraud).} TeMP-TraG (Gounoue et al., 2025) — пример метода для обнаружения мошенничества в транзакционных сетях: он вводит \emph{temporal message passing}, при котором недавние транзакции взвешиваются сильнее при агрегации сообщений【8†L113-L121】. Модель учитывает признаки ребер и мульти-графовую структуру. Эксперименты на реальных финансовых графах показывают существенное улучшение F1 по сравнению с классическими GNN【8†L39-L48】【8†L113-L121】. Концептуально это задача классификации узлов/ребер по признаку подозрительности. Обзор Cheng et al. (2024) охватывает более 100 работ GNN для финансового мошенничества и отмечает их высокую эффективность: GNN улавливают сложные реляционные паттерны и часто превосходят традиционные алгоритмы по всем метрикам【19†L60-L63】. Квантово-вдохновленные подходы (D’Amico et al., 2025) добавляют tensor-GNN и ансамбли (QBoost) и демонстрируют, что комплексный подход может достигать F2~0.75 при детектировании нелегальных транзакций【41†L49-L56】. 

\paragraph{Node classification и другие цели.} Ghimire et al. (2026) сравнивают евклидовы и гиперболические GNN для классификации узлов в Bitcoin-графе【36†L67-L75】【33†L166-L174】. Они показывают, что гиперболические эмбеддинги лучше захватывают иерархическую структуру многопрыжковых транзакций и дают более высокие метрики (Precision, Recall, F1) при больших глубинах модели【33†L166-L174】. Это важно для задач классификации типов участников и выявления сложных связных структур. Другие работы классифицируют узлы по активности (например, определяют «дневных» или «месячных» трейдеров по интервалам транзакций)【17†L859-L863】. Такие node-level предсказания дают временной взгляд на поведение отдельных аккаунтов.

\subsection{Сравнительная таблица работ}

\begin{table}[ht]
\caption{Сравнение ключевых работ по прогнозированию транзакций.}
\centering
\begin{tabular}{|p{2.3cm}|p{1.3cm}|p{2cm}|p{2cm}|p{2cm}|p{3cm}|p{2cm}|p{3cm}|}
\hline
Авторы (год) & arXiv ID & Домен & Тип графа & Цель(таргет) & Метод & Метрики & Ключевой вывод \\
\hline
Kandanaarachchi (2024) & 2401.04280 & Общая (биоинформатика и др.) & неряд., не взв., динамич. & Полная структура графа (новые узлы/рёбра) & TS на степенях + Flux Balance Analysis & (synthetic) точность структуры & Прогноз структуры будущего графа через линейн. оптимизацию; демонстрирует возможность учёта новых узлов【1†L43-L51】. \\
\hline
Takahashi et al. (2024) & 2409.08718 & Крипто/Банк & направ., взв., динамич. & Наличие рёбер и вес (поток) & Временные GNN (self-attention) + прогноз объёма & MAPE/MAE на весах & Разделяют задачу: сначала прогнозируют соотношение переводов, затем общий объём; умножение результатов даёт предсказание потока【6†L22-L30】. \\
\hline
Gounoue et al. (2025) & 2503.16901 & Финансовый & направ., взв., динамич. мульти & Классификация узлов/рёбер (AML) & Temporal Message Passing GNN (TeMP-TraG) & F1-score (fraud detection) & Учитывает временную близость транзакций в агрегации сообщений; улучшает SOTA GNN для обнаружения финансовых преступлений【8†L113-L121】. \\
\hline
Ghimire et al. (2026) & 2603.16080 & Bitcoin-сеть & направ., взв., статич. & Классификация узлов (типы адресов) & Евклидов и гипербол. GNN & Precision, Recall, F1 & Гиперболические GNN значительно превосходят евклидовы при увеличении глубины (multi-hop)【33†L166-L174】. \\
\hline
Wei et al. (2022) & 2007.07993 & Bitcoin & направ., взв., динамич. & Прогноз существования транзакций между адресами & Две временные графовые модели (доступность + паттерн суммы) + эмбеддинг & accuracy (Forecast) & Используют time-decaying reachability и amount-графы, обучают DNN; достигают ~60\% точности и +50\% к статическому бенчмарку【24†L25-L30】. \\
\hline
Wang et al. (2023) & 2310.01015 & Ethereum + off-chain & гетерогенный, динамич. & Link prediction \& детектирование мошенничества & Сбор и интеграция on-/off-chain данных, GNN-бенчмарк & AUC, F1 (phishing) & Создали большой датасет EX-Graph; показали, что добавление off-chain признаков улучшает прогноз транзакций и помогает в обнаружении wash-трейдинга. \\
\hline
Liang et al. (2024) & 2309.15728 & Общая & ненапр., взв., статич. & Прогноз веса ребра & Line-Graph CNN (LGLWP) на субграфе окрестности & RMSE, MAE на весе & Предлагают спускаться к line-графу, чтобы учить глубокие признаки ребра; существенно превосходят предыдущие методы по регрессии веса【22†L9-L17】. \\
\hline
Ma et al. (2025) & 2502.01029 & Bitcoin (fees) & статич. (тенденция) & Прогноз сборов транзакций (time series) & SARIMAX, Prophet, TFT и гибридные & RMSE, MAPE (fees) & Тщательное сравнение 6 моделей: традиционный SARIMAX показал лучшую точность, глубокие модели уступили из-за малого объёма данных【26†L61-L69】. \\
\hline
Cheng et al. (2024) & 2411.05815 & Финансы & финансовые сети & Обзор GNN для AML & Обзор литературы & — & Вывод: GNN отлично захватывают сложные взаимосвязи в финансовых графах и значительно превосходят классические методы детектирования мошенничества【19†L60-L63】. \\
\hline
D’Amico et al. (2025) & 2508.09237 & Blockchain/AML & направ., взв., статич. & Классификация узлов (AML) & Tensor GNN + QBoost ансамбль & F2-score (fraud) & Сочетание квантово-вдохновленного GNN (CP-декомпозиция) и ансамблей (QBoost) даёт F2≈0.748; указывает на потенциал квантовых методов в AML【41†L49-L56】. \\
\hline
\end{tabular}
\end{table}

\subsection{Обоснование выбора \emph{Link Prediction}}

Link prediction как формулировка оптимальна для нашей задачи по нескольким причинам. Во-первых, целевые метки (факт совершения транзакции между двумя узлами в будущем) естественно доступны: из исторических данных мы знаем, какие новые связи возникли и можем на них учиться【24†L25-L30】【17†L859-L863】. Во-вторых, для link prediction существуют устоявшиеся метрики качества (AUC, MRR), применимые к разреженным динамическим графам【15†L731-L739】【17†L859-L863】. Третье: аннотирование будущих ребер тривиально (вся необходимая информация уже в данных), тогда как альтернативные задачи сложнее размечать: вес транзакции (объём) может требовать точных данных о суммах и подвержен высоким шумам; аномалии нуждаются в метках мошенничества, которых мало и которые субъективны【19†L60-L63】. Кроме того, целевые задачи link prediction лучше согласуются с прогнозом сетевой активности: мы хотим предсказать структуру потока платежей, а не агрегированные суммы или отдельные события. В сравнении, задачи прогнозирования весов ребер часто фокусируются на уже существующих связях и решаются как регрессия【6†L22-L30】【22†L9-L17】, прогноз метрик транзакций (мемпул-фийсы) — на агрегированных временных рядах【26†L61-L69】, а обнаружение аномалий — на классификации узлов/ребер. Все эти альтернативы менее прямо решают задачу прогнозов новых транзакций. Поэтому именно link prediction отвечаете большинству критериев: метки есть и легкодоступны, метрики определены, сложность разметки минимальна, и постановка напрямую нацелена на прогнозирование графовой структуры транзакций. В связи с этим мы выбираем \emph{link prediction} как основную задачу в нашей работе.

\subsection*{Выводы}
В литературе показано, что link prediction является мощным и оправданным выбором для прогнозирования динамики графов транзакций. Модели GNN и представления графов хорошо справляются с задачей предсказания новых рёбер в активных криптографических сетях【17†L859-L863】【24†L25-L30】. На основе проведённого обзора мы приходим к выводу, что фокус на link prediction позволит использовать богатые структурные свойства графов, доступные цели и метрики, и даст наилучшую применимость для нашей постановки. В последующих разделах мы опишем выбранный нами подход к задаче link prediction на графе транзакций.

\begin{thebibliography}{99}
\bibitem{Kandanaarachchi2024} S. Kandanaarachchi, ``Predicting the structure of dynamic graphs,'' \textit{arXiv preprint} arXiv:2401.04280 (2024). 
\bibitem{Takahashi2024} S. Takahashi et al., ``Dynamic Link and Flow Prediction in Bank Transfer Networks,'' \textit{arXiv preprint} arXiv:2409.08718 (2024). 
\bibitem{Gounoue2025} S. Gounoue et al., ``TeMP-TraG: Edge-based Temporal Message Passing in Transaction Graphs,'' \textit{arXiv preprint} arXiv:2503.16901 (2025). 
\bibitem{Ghimire2026} A. Ghimire, S. Murad, N. Rahimi, ``A Depth-Aware Comparative Study of Euclidean and Hyperbolic Graph Neural Networks on Bitcoin Transaction Systems,'' \textit{arXiv preprint} arXiv:2603.16080 (2026). 
\bibitem{Wei2022} W. Wei, Q. Zhang, L. Liu, ``Bitcoin Transaction Forecasting with Deep Network Representation Learning,'' \textit{arXiv preprint} arXiv:2007.07993 (2022). 
\bibitem{Wang2023} Q. Wang et al., ``EX-Graph: A Pioneering Dataset Bridging Ethereum and X,'' \textit{arXiv preprint} arXiv:2310.01015 (2023). 
\bibitem{Liang2024} J. Liang et al., ``Line Graph Neural Networks for Link Weight Prediction,'' \textit{arXiv preprint} arXiv:2309.15728 (2023). 
\bibitem{Ma2025} J. Ma, E. Mahmoudinia, ``Comprehensive Modeling Approaches for Forecasting Bitcoin Transaction Fees: A Comparative Study,'' \textit{arXiv preprint} arXiv:2502.01029 (2025). 
\bibitem{Cheng2024} D. Cheng et al., ``Graph Neural Networks for Financial Fraud Detection: A Review,'' \textit{arXiv preprint} arXiv:2411.05815 (2024). 
\bibitem{Damico2025} L. D'Amico et al., ``Blockchain Network Analysis using Quantum Inspired Graph Neural Networks \& Ensemble Models,'' \textit{arXiv preprint} arXiv:2508.09237 (2025). 
\end{thebibliography}
